% Figure: Left-ventricular pressure-volume loop with the four phases
% (filling, isovolumetric contraction, ejection, isovolumetric
% relaxation) annotated.
\begin{figure}[htbp]
\centering
\begin{tikzpicture}[
  axis/.style={-{Stealth}, thick, draw=engblack},
  loop/.style={very thick, engred},
  edge/.style={thin, engblack},
  annot/.style={font=\scriptsize, align=left},
]

\draw[axis] (0, 0) -- (10, 0) node[right, font=\small] {LV volume (mL)};
\draw[axis] (0, 0) -- (0, 7) node[above, font=\small] {LV pressure (mmHg)};

\foreach \x/\xt in {2/40, 4/80, 6/120, 8/160} {
  \draw (\x, 0) -- (\x, -0.1) node[below, font=\scriptsize] {\xt};
}
\foreach \y/\yt in {1/20, 2.5/50, 4/80, 5.5/110, 6.5/140} {
  \draw (0, \y) -- (-0.1, \y) node[left, font=\scriptsize] {\yt};
}

% PV loop: simple parallelogram-with-curves approximation.
% Vertices: A = filling start (50 mL, 5 mmHg)
%           B = end-diastole / mitral close (120 mL, 8 mmHg)
%           C = aortic open (120 mL, 80 mmHg)
%           D = end-systole / aortic close (50 mL, 100 mmHg)
% A: (2.5, 0.25), B: (6, 0.4), C: (6, 4), D: (2.5, 5)

% Filling (A->B): along EDPVR (passive end-diastolic curve)
\draw[loop] (2.5, 0.25) .. controls (4.5, 0.3) and (5.5, 0.35) .. (6, 0.4);
% Isovolumetric contraction (B->C)
\draw[loop] (6, 0.4) -- (6, 4);
% Ejection (C->D): along the high-pressure plateau, slight curve
\draw[loop] (6, 4) .. controls (5.0, 4.8) and (3.5, 5.1) .. (2.5, 5);
% Isovolumetric relaxation (D->A)
\draw[loop] (2.5, 5) -- (2.5, 0.25);

% ESPVR (end-systolic pressure-volume relation): linear, dashed
\draw[edge, dashed] (1.0, 0) -- (4.0, 7);
\node[annot, anchor=west] at (3.5, 6.5) {ESPVR};

% EDPVR: low and nearly flat
\draw[edge, dashed] (1.0, 0) .. controls (4.0, 0.1) and (6.5, 0.5) .. (8.0, 1.7);
\node[annot, anchor=west] at (7.5, 1.3) {EDPVR};

% Vertex labels
\fill[engred] (2.5, 0.25) circle (2pt);
\node[annot, anchor=east] at (2.4, 0.4) {A (MV open)};
\fill[engred] (6, 0.4) circle (2pt);
\node[annot, anchor=west] at (6.1, 0.3) {B (MV close)};
\fill[engred] (6, 4) circle (2pt);
\node[annot, anchor=west] at (6.1, 4.0) {C (AV open)};
\fill[engred] (2.5, 5) circle (2pt);
\node[annot, anchor=east] at (2.4, 5.2) {D (AV close)};

% Stroke volume and stroke work
\node[font=\small, engred] at (4.3, 2.7) {SV};
\node[annot] at (4.3, 2.3) {stroke work\\$\approx 1\,\text{J}/\text{beat}$};

% Phase annotations
\node[annot] at (7.5, 0.1) {filling (A$\to$B)};
\node[annot] at (6.5, 2.5) {iso. contr.\\(B$\to$C)};
\node[annot] at (4.3, 4.8) {ejection (C$\to$D)};
\node[annot] at (1.0, 2.5) {iso.\\relax.\\(D$\to$A)};

\node[font=\bfseries, anchor=south] at (5, 7.0) {Left ventricular pressure-volume loop};

\end{tikzpicture}
\caption{Left-ventricular pressure-volume loop, one cardiac cycle.
Phase A$\to$B is diastolic filling, governed by the end-diastolic
pressure-volume relation (EDPVR, lower dashed curve). Phase B$\to$C
is isovolumetric contraction (mitral valve closed, aortic valve not
yet open). Phase C$\to$D is ejection. Phase D$\to$A is isovolumetric
relaxation. End-systolic points lie on the end-systolic
pressure-volume relation (ESPVR, upper dashed line), whose slope is
the load-independent measure of contractility. The area inside the
loop is the stroke work per beat, $\approx 1\,\text{J}$ at rest. The
loop's geometry is the clinical fingerprint of cardiac function:
dilation moves it rightward, concentric hypertrophy upward and
leftward, failing hearts collapse the stroke volume.}
\label{fig:vol06ch06:pv-loop}
\end{figure}
